\chapter{Radioaktivitaet und Kernzerfaelle}

\section{Ziel dieses Kapitels}

Dieses Kapitel behandelt radioaktive Zerfaelle von Atomkernen.
Im Mittelpunkt stehen das Zerfallsgesetz, Halbwertszeit, Aktivitaet,
Alpha-Zerfall, Beta-Zerfall, Gamma-Zerfall, innere Konversion und die wichtigsten
Erhaltungssaetze.

\begin{notebox}
Dieses Kapitel behandelt die Stichwoerter:
\begin{itemize}
    \item radioaktiver Zerfall,
    \item Zerfallsgesetz,
    \item Zerfallskonstante,
    \item Halbwertszeit,
    \item Aktivitaet,
    \item Zerfallsreihen,
    \item Alpha-Zerfall,
    \item Geiger-Nuttall-Regel,
    \item Beta-Zerfall,
    \item Neutrino,
    \item Elektroneneinfang,
    \item Gamma-Zerfall,
    \item innere Konversion,
    \item Multipolstrahlung,
    \item Erhaltungssaetze.
\end{itemize}
\end{notebox}

\section{Was bedeutet radioaktiver Zerfall?}

Radioaktive Kerne sind instabil.
Sie wandeln sich spontan in andere Kerne oder andere Energiezustaende um.

\begin{definitionbox}
Ein \emph{radioaktiver Zerfall} ist eine spontane Umwandlung eines instabilen
Atomkerns.

Dabei koennen Teilchen und Strahlung ausgesendet werden.
\end{definitionbox}

\begin{conceptbox}
Radioaktiver Zerfall ist ein statistischer Prozess.
Man kann fuer einen einzelnen Kern nicht vorhersagen, wann er zerfaellt.

Man kann aber fuer eine grosse Anzahl gleicher Kerne sehr genau vorhersagen,
wie viele Kerne pro Zeitintervall im Mittel zerfallen.
\end{conceptbox}

\begin{examtipbox}
Wichtig:
Radioaktiver Zerfall ist spontan und statistisch.
Die Halbwertszeit beschreibt nicht den Zerfallszeitpunkt eines einzelnen Kerns,
sondern das Verhalten einer grossen Anzahl gleicher Kerne.
\end{examtipbox}

\section{Zerfallsgesetz}

Sei $N(t)$ die Anzahl der noch nicht zerfallenen Kerne zur Zeit $t$.

\begin{definitionbox}
Die \emph{Zerfallskonstante} $\lambda$ gibt die Zerfallswahrscheinlichkeit pro Zeit
fuer einen einzelnen Kern an.
\end{definitionbox}

\begin{formulabox}
Das Zerfallsgesetz lautet:
\[
    N(t)=N_0 e^{-\lambda t}.
\]
Dabei ist:
\begin{itemize}
    \item $N_0=N(0)$ die Anfangszahl der Kerne,
    \item $\lambda$ die Zerfallskonstante,
    \item $t$ die Zeit.
\end{itemize}
\end{formulabox}

\begin{derivationbox}
Die Anzahl der Zerfaelle pro Zeit ist proportional zur Anzahl der noch vorhandenen
Kerne:
\[
    \frac{\dd N}{\dd t}=-\lambda N.
\]
Trennung der Variablen ergibt:
\[
    \frac{\dd N}{N}=-\lambda \dd t.
\]
Integration liefert:
\[
    \ln N = -\lambda t + C.
\]
Mit $N(0)=N_0$ folgt:
\[
    N(t)=N_0 e^{-\lambda t}.
\]
\end{derivationbox}

\begin{conceptbox}
Das Minuszeichen in
\[
    \frac{\dd N}{\dd t}=-\lambda N
\]
bedeutet, dass die Zahl der noch nicht zerfallenen Kerne abnimmt.
\end{conceptbox}

\section{Halbwertszeit}

\begin{definitionbox}
Die \emph{Halbwertszeit} $T_{1/2}$ ist die Zeit, nach der nur noch die Haelfte der
urspruenglichen Kerne vorhanden ist.
\end{definitionbox}

\begin{formulabox}
\[
    N(T_{1/2})=\frac{N_0}{2}.
\]
Mit dem Zerfallsgesetz folgt:
\[
    \frac{N_0}{2}=N_0 e^{-\lambda T_{1/2}}.
\]
Daraus:
\[
    T_{1/2}=\frac{\ln 2}{\lambda}.
\]
\end{formulabox}

\begin{derivationbox}
Aus
\[
    \frac{1}{2}=e^{-\lambda T_{1/2}}
\]
folgt durch Logarithmieren:
\[
    \ln\left(\frac{1}{2}\right)=-\lambda T_{1/2}.
\]
Da
\[
    \ln\left(\frac{1}{2}\right)=-\ln 2
\]
ist, erhaelt man:
\[
    T_{1/2}=\frac{\ln 2}{\lambda}.
\]
\end{derivationbox}

\begin{examtipbox}
Wenn die Halbwertszeit bekannt ist, kann man die Zerfallskonstante bestimmen:
\[
    \lambda=\frac{\ln 2}{T_{1/2}}.
\]
\end{examtipbox}

\section{Aktivitaet}

\begin{definitionbox}
Die \emph{Aktivitaet} $A(t)$ ist die Anzahl der Zerfaelle pro Zeit.
\end{definitionbox}

\begin{formulabox}
\[
    A(t)=-\frac{\dd N}{\dd t}.
\]
Mit dem Zerfallsgesetz:
\[
    A(t)=\lambda N(t).
\]
Also:
\[
    A(t)=A_0 e^{-\lambda t},
    \qquad
    A_0=\lambda N_0.
\]
\end{formulabox}

\begin{definitionbox}
Die SI-Einheit der Aktivitaet ist das Becquerel:
\[
    1\,\mathrm{Bq}=1\,\mathrm{s}^{-1}.
\]
Ein Becquerel bedeutet ein Zerfall pro Sekunde.
\end{definitionbox}

\begin{conceptbox}
Aktivitaet ist nicht dasselbe wie Anzahl der Kerne.
Eine Probe kann viele Kerne enthalten, aber eine kleine Aktivitaet haben, wenn die
Halbwertszeit sehr gross ist.
\end{conceptbox}

\begin{examtipbox}
Typischer Rechenweg:
\begin{enumerate}
    \item Aus $T_{1/2}$ die Zerfallskonstante bestimmen:
    \[
        \lambda=\frac{\ln 2}{T_{1/2}}.
    \]
    \item Dann Aktivitaet berechnen:
    \[
        A=\lambda N.
    \]
\end{enumerate}
\end{examtipbox}

\section{Mittlere Lebensdauer}

\begin{definitionbox}
Die \emph{mittlere Lebensdauer} $\tau$ ist der Erwartungswert der Lebensdauer eines
radioaktiven Kerns.
\end{definitionbox}

\begin{formulabox}
\[
    \tau=\frac{1}{\lambda}.
\]
Der Zusammenhang mit der Halbwertszeit ist:
\[
    T_{1/2}=\tau\ln 2.
\]
\end{formulabox}

\begin{conceptbox}
Die Halbwertszeit $T_{1/2}$ ist die Zeit, bis die Haelfte einer grossen Menge
zerfallen ist.

Die mittlere Lebensdauer $\tau$ ist der Mittelwert der individuellen Zerfallszeiten.
\end{conceptbox}

\section{Zerfallsreihen}

Schwere Kerne zerfallen oft nicht in einem einzigen Schritt in einen stabilen Kern.
Stattdessen entsteht eine Zerfallsreihe.

\begin{definitionbox}
Eine \emph{Zerfallsreihe} ist eine Folge radioaktiver Zerfaelle, bei der das
Tochternuklid selbst wieder radioaktiv sein kann.
\end{definitionbox}

\begin{conceptbox}
In natuerlichen Zerfallsreihen treten haeufig Alpha- und Beta-Zerfaelle auf.
Jeder Zerfall veraendert $A$ und $Z$ nach festen Regeln.

Alpha-Zerfall verringert:
\[
    A \rightarrow A-4,
    \qquad
    Z \rightarrow Z-2.
\]

Beta-Minus-Zerfall laesst $A$ gleich und erhoeht:
\[
    Z \rightarrow Z+1.
\]
\end{conceptbox}

\begin{notebox}
Die Vorlesungsabbildung zeigt natuerliche radioaktive Zerfallsreihen, zum Beispiel
Uran-Radium-Reihe, Uran-Actinium-Reihe, Thorium-Reihe und Neptunium-Reihe.
In solchen Darstellungen sieht man die charakteristischen Schritte von Alpha- und
Beta-Zerfaellen im $(A,Z)$-Diagramm.
\end{notebox}

\begin{examtipbox}
Bei Zerfallsreihen ist wichtig:
\begin{itemize}
    \item Alpha-Zerfall verschiebt den Kern um $\Delta A=-4$ und $\Delta Z=-2$.
    \item Beta-Minus-Zerfall verschiebt den Kern um $\Delta A=0$ und $\Delta Z=+1$.
    \item Beta-Plus-Zerfall oder Elektroneneinfang verschiebt um $\Delta A=0$ und $\Delta Z=-1$.
\end{itemize}
\end{examtipbox}

\section{Alpha-Zerfall}

Beim Alpha-Zerfall sendet ein instabiler Kern ein Alpha-Teilchen aus.
Ein Alpha-Teilchen ist ein Helium-4-Kern.

\begin{definitionbox}
Ein \emph{Alpha-Teilchen} ist ein ${}^{4}_{2}\mathrm{He}$-Kern.
Es besteht aus zwei Protonen und zwei Neutronen.
\end{definitionbox}

\begin{formulabox}
Ein allgemeiner Alpha-Zerfall hat die Form:
\[
    {}^{A}_{Z}X
    \rightarrow
    {}^{A-4}_{Z-2}Y
    +
    {}^{4}_{2}\mathrm{He}.
\]
\end{formulabox}

\begin{conceptbox}
Beim Alpha-Zerfall gelten:
\[
    A \rightarrow A-4,
    \qquad
    Z \rightarrow Z-2.
\]
Die Massenzahl nimmt also um $4$ ab, die Protonenzahl um $2$.
\end{conceptbox}

\begin{examplebox}
Ein Beispiel ist:
\[
    {}^{238}_{92}\mathrm{U}
    \rightarrow
    {}^{234}_{90}\mathrm{Th}
    +
    {}^{4}_{2}\mathrm{He}.
\]
Hier gilt:
\[
    238 \rightarrow 234+4,
\]
\[
    92 \rightarrow 90+2.
\]
Massenzahl und Ladung sind erhalten.
\end{examplebox}

\section{Energiebilanz beim Alpha-Zerfall}

\begin{definitionbox}
Der $Q$-Wert des Alpha-Zerfalls ist die freiwerdende Energie:
\[
    Q_\alpha
    =
    \left(
        m_{\mathrm{Mutter}}
        -
        m_{\mathrm{Tochter}}
        -
        m_{\alpha}
    \right)c^2.
\]
\end{definitionbox}

\begin{conceptbox}
Alpha-Zerfall ist energetisch moeglich, wenn
\[
    Q_\alpha>0
\]
gilt.

Die freiwerdende Energie erscheint vor allem als kinetische Energie des Alpha-Teilchens
und des Tochterkerns.
Da der Tochterkern viel schwerer ist, bekommt das Alpha-Teilchen den groessten Anteil
der kinetischen Energie.
\end{conceptbox}

\begin{formulabox}
Impulserhaltung verlangt:
\[
    \vec p_\alpha + \vec p_{\mathrm{Tochter}} = 0.
\]
Daher haben Alpha-Teilchen und Tochterkern betragsgleiche Impulse, aber entgegengesetzte
Richtungen.
\end{formulabox}

\section{Alpha-Zerfall als Tunneleffekt}

Klassisch waere der Alpha-Zerfall schwer zu verstehen.
Das Alpha-Teilchen ist im Kern gebunden, sieht aber ausserhalb des Kerns eine
Coulomb-Barriere.

\begin{conceptbox}
Der Alpha-Zerfall ist ein quantenmechanischer Tunneleffekt.

Das Alpha-Teilchen besitzt nicht genug Energie, um die Coulomb-Barriere klassisch
zu ueberwinden.
Quantenmechanisch kann es aber mit endlicher Wahrscheinlichkeit durch die Barriere
tunneln.
\end{conceptbox}

\begin{notebox}
Die Vorlesungsabbildung zum Alpha-Zerfall zeigt ein Potential mit gebundenen
Zustaenden im Kerninneren und einer Coulomb-Barriere ausserhalb.
Das Alpha-Teilchen kann durch diese Barriere tunneln.
\end{notebox}

\begin{examtipbox}
Wichtige Aussage:
Die sehr unterschiedlichen Halbwertszeiten von Alpha-Strahlern entstehen durch die
starke Energieabhaengigkeit der Tunnelwahrscheinlichkeit.
Schon kleine Unterschiede in der Alpha-Energie koennen grosse Unterschiede in der
Halbwertszeit bewirken.
\end{examtipbox}

\section{Geiger-Nuttall-Regel}

Die Geiger-Nuttall-Regel beschreibt den Zusammenhang zwischen Alpha-Energie und
Halbwertszeit.

\begin{definitionbox}
Die \emph{Geiger-Nuttall-Regel} ist ein empirischer Zusammenhang zwischen der
Reichweite beziehungsweise Energie von Alpha-Teilchen und der Zerfallskonstante
beziehungsweise Halbwertszeit des Alpha-Strahlers.
\end{definitionbox}

\begin{conceptbox}
Alpha-Teilchen mit groesserer Energie koennen die Coulomb-Barriere leichter
durchtunneln.

Daher gilt qualitativ:
\[
    E_\alpha \text{ gross}
    \quad\Rightarrow\quad
    T_{1/2} \text{ klein}.
\]
\[
    E_\alpha \text{ klein}
    \quad\Rightarrow\quad
    T_{1/2} \text{ gross}.
\]
\end{conceptbox}

\begin{notebox}
In den Vorlesungsfolien ist eine Geiger-Nuttall-Kurve dargestellt.
Sie zeigt experimentelle Werte fuer Alpha-Strahler aus natuerlichen Zerfallsreihen.
\end{notebox}

\section{Reichweite und Ionisation von Alpha-Teilchen}

Alpha-Teilchen sind schwer und zweifach positiv geladen.
Sie ionisieren Materie stark.

\begin{conceptbox}
Beim Durchgang durch Materie verliert ein Alpha-Teilchen Energie durch Ionisation
und Anregung von Atomen.

Die spezifische Ionisation steigt gegen Ende der Bahn stark an.
Dieses Maximum nennt man Bragg-Peak.
\end{conceptbox}

\begin{definitionbox}
Der \emph{Bragg-Peak} ist das Maximum der Energieabgabe eines geladenen Teilchens
kurz vor dem Ende seiner Reichweite in Materie.
\end{definitionbox}

\begin{notebox}
Die Vorlesungsabbildung zur spezifischen Ionisation von Alpha-Teilchen zeigt,
dass die Ionisation entlang der Bahn ansteigt und kurz vor dem Stillstand ein Maximum
erreicht.
\end{notebox}

\section{Beta-Zerfall}

Beim Beta-Zerfall wird ein Neutron in ein Proton oder ein Proton in ein Neutron
umgewandelt.
Dabei wirken die schwache Wechselwirkung und Neutrinos.

\begin{definitionbox}
Beim \emph{Beta-Zerfall} aendert sich die Protonenzahl $Z$, aber die Massenzahl $A$
bleibt gleich.
\end{definitionbox}

\subsection{Beta-Minus-Zerfall}

\begin{formulabox}
Beim $\beta^-$-Zerfall gilt:
\[
    n \rightarrow p + e^- + \overline{\nu}_e.
\]
Fuer den Kern:
\[
    {}^{A}_{Z}X
    \rightarrow
    {}^{A}_{Z+1}Y
    +
    e^-
    +
    \overline{\nu}_e.
\]
\end{formulabox}

\begin{conceptbox}
Beim $\beta^-$-Zerfall steigt die Protonenzahl:
\[
    Z \rightarrow Z+1.
\]
Die Massenzahl bleibt gleich:
\[
    A \rightarrow A.
\]
\end{conceptbox}

\subsection{Beta-Plus-Zerfall}

\begin{formulabox}
Beim $\beta^+$-Zerfall gilt:
\[
    p \rightarrow n + e^+ + \nu_e.
\]
Fuer den Kern:
\[
    {}^{A}_{Z}X
    \rightarrow
    {}^{A}_{Z-1}Y
    +
    e^+
    +
    \nu_e.
\]
\end{formulabox}

\begin{conceptbox}
Beim $\beta^+$-Zerfall sinkt die Protonenzahl:
\[
    Z \rightarrow Z-1.
\]
Die Massenzahl bleibt gleich:
\[
    A \rightarrow A.
\]
\end{conceptbox}

\subsection{Elektroneneinfang}

\begin{definitionbox}
Beim \emph{Elektroneneinfang} faengt der Kern ein Huellenelektron ein.
Ein Proton wird dadurch in ein Neutron umgewandelt.
\end{definitionbox}

\begin{formulabox}
\[
    p + e^- \rightarrow n + \nu_e.
\]
Fuer den Kern:
\[
    {}^{A}_{Z}X + e^-
    \rightarrow
    {}^{A}_{Z-1}Y + \nu_e.
\]
\end{formulabox}

\begin{conceptbox}
Elektroneneinfang und $\beta^+$-Zerfall haben im Kern denselben Effekt:
\[
    Z \rightarrow Z-1,
    \qquad
    A \rightarrow A.
\]
\end{conceptbox}

\section{Beta-Zerfall und Neutrino}

Das Beta-Spektrum ist kontinuierlich.
Das war historisch ein Hinweis darauf, dass neben Elektron oder Positron noch ein
weiteres Teilchen beteiligt ist: das Neutrino.

\begin{conceptbox}
Bei einem Zweikoerperzerfall waere die Energie des ausgesendeten Elektrons fest.

Beim Beta-Zerfall beobachtet man aber ein kontinuierliches Energiespektrum.
Das bedeutet:
Die Zerfallsenergie wird auf mehrere Teilchen verteilt.

Das Neutrino beziehungsweise Antineutrino traegt einen variablen Anteil der Energie
und des Impulses.
\end{conceptbox}

\begin{notebox}
Die Vorlesungsfolie mit Bezug auf KATRIN gehoert zum Beta-Zerfall.
KATRIN untersucht das Endpunktspektrum des Tritium-Beta-Zerfalls, um Informationen
ueber die Neutrinomasse zu gewinnen.
\end{notebox}

\begin{examtipbox}
Warum braucht man das Neutrino beim Beta-Zerfall?

Damit Energie, Impuls und Drehimpuls erhalten bleiben und das kontinuierliche
Beta-Spektrum erklaert wird.
\end{examtipbox}

\section{Erhaltungssaetze beim Beta-Zerfall}

Beta-Zerfaelle muessen mehrere Erhaltungssaetze erfuellen.

\begin{notebox}
Wichtige Erhaltungsgroessen sind:
\begin{itemize}
    \item elektrische Ladung,
    \item Baryonenzahl,
    \item Leptonenzahl,
    \item Energie,
    \item Impuls,
    \item Drehimpuls.
\end{itemize}
\end{notebox}

\begin{examplebox}
Beim Neutronzerfall:
\[
    n \rightarrow p + e^- + \overline{\nu}_e.
\]

Ladung:
\[
    0 = (+1) + (-1) + 0.
\]

Baryonenzahl:
\[
    1 = 1 + 0 + 0.
\]

Leptonenzahl:
\[
    0 = 0 + 1 - 1.
\]
\end{examplebox}

\begin{examtipbox}
Bei Beta-Zerfaellen ist die Leptonenzahl besonders wichtig.

Ein Elektron hat
\[
    L_e=+1.
\]
Ein Elektron-Antineutrino hat
\[
    L_e=-1.
\]
Daher ist beim $\beta^-$-Zerfall die Leptonenzahl erhalten.
\end{examtipbox}

\section{Gamma-Zerfall}

Beim Gamma-Zerfall aendert sich nicht die Zusammensetzung des Kerns, sondern nur
sein Energiezustand.

\begin{definitionbox}
Beim \emph{Gamma-Zerfall} geht ein angeregter Kern in einen Zustand niedrigerer
Energie ueber und sendet dabei ein Photon aus.
\end{definitionbox}

\begin{formulabox}
\[
    {}^{A}_{Z}X^*
    \rightarrow
    {}^{A}_{Z}X
    +
    \gamma.
\]
\end{formulabox}

\begin{conceptbox}
Beim Gamma-Zerfall bleiben Massenzahl und Protonenzahl gleich:
\[
    A \rightarrow A,
    \qquad
    Z \rightarrow Z.
\]
Nur die innere Energie des Kerns aendert sich.
\end{conceptbox}

\begin{definitionbox}
Die Energie des Gamma-Photons ist naeherungsweise die Energiedifferenz zwischen
Anfangs- und Endzustand:
\[
    E_\gamma \approx E_i-E_f.
\]
\end{definitionbox}

\section{Gamma-Uebergaenge und Multipolstrahlung}

Gamma-Uebergaenge unterliegen Auswahlregeln.
Diese haengen mit Drehimpuls und Paritaet zusammen.

\begin{definitionbox}
Die \emph{Multipolordnung} $l$ beschreibt den Drehimpuls, der durch das emittierte
Photon getragen wird.
\end{definitionbox}

\begin{notebox}
In den Vorlesungsfolien werden Einteilchen-Gamma-Uebergangswahrscheinlichkeiten fuer
verschiedene Multipolstrahlungen gezeigt:
\begin{itemize}
    \item $l=1$: Dipolstrahlung,
    \item $l=2$: Quadrupolstrahlung,
    \item $l=3$: Oktopolstrahlung.
\end{itemize}

Ausserdem wird zwischen elektrischen Uebergaengen $E l$ und magnetischen Uebergaengen
$M l$ unterschieden.
\end{notebox}

\begin{conceptbox}
Je hoeher die Multipolordnung ist, desto unwahrscheinlicher ist der Uebergang im
Allgemeinen.

Deshalb koennen bestimmte angeregte Kernzustaende relativ langlebig sein.
Solche langlebigen angeregten Zustaende nennt man Kernisomere.
\end{conceptbox}

\begin{examtipbox}
Wichtige Begriffe:
\[
    E1, \quad M1, \quad E2, \quad M2.
\]
Dabei steht $E$ fuer elektrischen Uebergang und $M$ fuer magnetischen Uebergang.
Die Zahl gibt die Multipolordnung an.
\end{examtipbox}

\section{Innere Konversion}

Gamma-Abregung ist nicht die einzige Moeglichkeit, wie ein angeregter Kern Energie
abgeben kann.

\begin{definitionbox}
Bei der \emph{inneren Konversion} uebertraegt ein angeregter Kern seine Energie direkt
auf ein Huellenelektron.
Dieses Elektron wird aus dem Atom herausgeloest.
\end{definitionbox}

\begin{conceptbox}
Innere Konversion ist kein Beta-Zerfall.

Beim Beta-Zerfall entsteht ein Elektron im Kernprozess.
Bei innerer Konversion war das Elektron bereits in der Atomhülle vorhanden und wird
durch die Kernabregung herausgeschlagen.
\end{conceptbox}

\begin{formulabox}
Die kinetische Energie des Konversionselektrons ist naeherungsweise:
\[
    E_{\mathrm{kin}}
    =
    E_{\mathrm{Uebergang}}
    -
    E_{\mathrm{Bindung}}
\]
wobei $E_{\mathrm{Bindung}}$ die Bindungsenergie des Huellenelektrons ist.
\end{formulabox}

\begin{examtipbox}
Unterschied merken:

Gamma-Zerfall:
\[
    \text{Kernenergie} \rightarrow \gamma\text{-Photon}.
\]

Innere Konversion:
\[
    \text{Kernenergie} \rightarrow \text{kinetische Energie eines Huellenelektrons}.
\]
\end{examtipbox}

\section{Vergleich der Zerfallsarten}

\begin{center}
\begin{tabular}{|l|l|l|l|}
\hline
Zerfall & Aenderung von $A$ & Aenderung von $Z$ & ausgesendet \\
\hline
$\alpha$ & $A-4$ & $Z-2$ & ${}^{4}_{2}\mathrm{He}$ \\
\hline
$\beta^-$ & $A$ & $Z+1$ & $e^-+\overline{\nu}_e$ \\
\hline
$\beta^+$ & $A$ & $Z-1$ & $e^+ + \nu_e$ \\
\hline
Elektroneneinfang & $A$ & $Z-1$ & $\nu_e$ \\
\hline
$\gamma$ & $A$ & $Z$ & $\gamma$ \\
\hline
innere Konversion & $A$ & $Z$ & Huellenelektron \\
\hline
\end{tabular}
\end{center}

\begin{examtipbox}
Die Tabelle ist sehr pruefungsrelevant.
Bei Zerfallsaufgaben zuerst immer klaeren:
\begin{itemize}
    \item Was passiert mit $A$?
    \item Was passiert mit $Z$?
    \item Welche Teilchen werden ausgesendet?
    \item Welche Erhaltungssaetze muessen gelten?
\end{itemize}
\end{examtipbox}

\section{Mini-Zusammenfassung}

\begin{notebox}
Die wichtigsten Punkte dieses Kapitels sind:
\begin{itemize}
    \item Radioaktiver Zerfall ist statistisch.
    \item Das Zerfallsgesetz lautet:
    \[
        N(t)=N_0e^{-\lambda t}.
    \]
    \item Die Halbwertszeit ist:
    \[
        T_{1/2}=\frac{\ln 2}{\lambda}.
    \]
    \item Die Aktivitaet ist:
    \[
        A(t)=\lambda N(t).
    \]
    \item Alpha-Zerfall:
    \[
        A\rightarrow A-4,
        \qquad
        Z\rightarrow Z-2.
    \]
    \item Beta-Minus-Zerfall:
    \[
        A\rightarrow A,
        \qquad
        Z\rightarrow Z+1.
    \]
    \item Beta-Plus-Zerfall und Elektroneneinfang:
    \[
        A\rightarrow A,
        \qquad
        Z\rightarrow Z-1.
    \]
    \item Gamma-Zerfall aendert nur den Energiezustand des Kerns.
    \item Alpha-Zerfall ist ein Tunneleffekt.
    \item Das kontinuierliche Beta-Spektrum fuehrt zum Neutrino.
    \item Gamma-Uebergaenge werden durch Multipolstrahlung beschrieben.
    \item Innere Konversion ist eine Alternative zur Gamma-Emission.
\end{itemize}
\end{notebox}

\section{Pruefungscheck}

\begin{examtipbox}
Nach diesem Kapitel solltest du folgende Fragen beantworten koennen:
\begin{enumerate}
    \item Was bedeutet radioaktiver Zerfall?
    \item Warum ist radioaktiver Zerfall statistisch?
    \item Wie lautet das Zerfallsgesetz?
    \item Wie haengen Halbwertszeit und Zerfallskonstante zusammen?
    \item Was ist Aktivitaet?
    \item Was ist eine Zerfallsreihe?
    \item Wie veraendern sich $A$ und $Z$ beim Alpha-Zerfall?
    \item Warum ist Alpha-Zerfall ein Tunneleffekt?
    \item Was beschreibt die Geiger-Nuttall-Regel?
    \item Was ist der Bragg-Peak?
    \item Wie veraendern sich $A$ und $Z$ beim $\beta^-$-Zerfall?
    \item Wie veraendern sich $A$ und $Z$ beim $\beta^+$-Zerfall?
    \item Was passiert beim Elektroneneinfang?
    \item Warum braucht man das Neutrino beim Beta-Zerfall?
    \item Welche Erhaltungssaetze gelten beim Beta-Zerfall?
    \item Was passiert beim Gamma-Zerfall?
    \item Was bedeutet Multipolstrahlung?
    \item Was ist innere Konversion?
    \item Wie unterscheidet man Gamma-Zerfall und innere Konversion?
\end{enumerate}
\end{examtipbox}